\chapter{Riduzione della Dimensionalità e Modelli Generativi}
\label{ch:modelli_generativi}

\section{Principal Component Analysis (PCA): teoria e limiti lineari}
\label{sec:pca}

La \emph{Principal Component Analysis} (PCA) rappresenta lo standard aureo classico per la riduzione della dimensionalità e l'estrazione di feature in ambito statistico e finanziario. Dal punto di vista algebrico, la PCA opera una trasformazione ortogonale lineare che proietta i dati originali in un nuovo sistema di coordinate, in cui gli assi (le componenti principali) sono allineati lungo le direzioni di massima varianza del dataset.

Considerando il vettore linearizzato di input $\mathbf{x} \in \mathbb{R}^D$ (che nel nostro dominio rappresenta gli elementi della matrice triangolare inferiore), la PCA ricerca una matrice di proiezione $\mathbf{W} \in \mathbb{R}^{D \times k}$ con $k < D$, tale per cui le colonne di $\mathbf{W}$ siano ortonormali ($\mathbf{W}^T \mathbf{W} = \mathbf{I}_k$). La proiezione nello spazio latente lineare è data da:
\begin{equation}
\mathbf{z} = \mathbf{W}^T \mathbf{x}
\end{equation}
dove $\mathbf{z} \in \mathbb{R}^k$ rappresenta il vettore delle coordinate compresse. La ricostruzione nello spazio originale si ottiene applicando la trasformazione inversa:
\begin{equation}
\mathbf{\hat{x}} = \mathbf{W} \mathbf{z} = \mathbf{W} \mathbf{W}^T \mathbf{x}
\end{equation}

L'ottimizzazione della PCA consiste nel trovare la matrice $\mathbf{W}$ che minimizzi l'errore quadratico medio di ricostruzione (MSE) globale, equivalente a massimizzare la traccia della matrice di covarianza proiettata. La soluzione analitica a questo problema coincide con l'autodecomposizione della matrice di covarianza dei dati: le colonne di $\mathbf{W}$ sono gli autovettori associati ai $k$ autovalori dominanti.

Tuttavia, l'applicazione della PCA per la modellizzazione di matrici di correlazione finanziarie incontra un limite insormontabile: la linearità intrinseca. La PCA assume che la \emph{manifold} (la varietà geometrica sotterranea) su cui risiedono i dati reali sia un iperpiano lineare. I mercati finanziari, al contrario, sono caratterizzati da dinamiche non lineari estreme, come asimmetrie nei co-movimenti sistemici durante le crisi (\emph{tail dependence}) e complesse interazioni gerarchiche di settore. Una trasformazione puramente lineare come la PCA è matematicamente incapace di catturare queste curvature topologiche, producendo ricostruzioni storiche che tendono ad appiattire l'informazione complessa.

\section{Autoencoder Lineari (Linear AE): ponte verso le reti neurali}
\label{sec:linear_ae}

Per superare il paradigma dell'algebra lineare classica, l'architettura dei modelli di \emph{deep learning} introduce gli \emph{Autoencoder} (AE). Un autoencoder è una rete neurale \emph{feed-forward} addestrata in modo non supervisionato con l'obiettivo di copiare il proprio input nel proprio output, passando attraverso uno strato intermedio compresso detto "collo di bottiglia" (\emph{bottleneck}).

L'architettura minima è composta da due funzioni: un \emph{Encoder} $f_{\phi}$ che mappa l'input nello spazio latente $\mathbf{z} = f_{\phi}(\mathbf{x})$, e un \emph{Decoder} $g_{\theta}$ che ricostruisce l'input $\mathbf{\hat{x}} = g_{\theta}(\mathbf{z})$.

Il punto di giunzione teorico tra reti neurali e algebra statistica è rappresentato dall'Autoencoder Lineare. Affinché tale modello replichi fedelmente la PCA, è necessario imporre vincoli architetturali rigorosi: le funzioni di attivazione devono essere strettamente lineari (identità), la rete deve essere priva di termini di \emph{bias} e non deve essere applicata alcuna forma di regolarizzazione sui pesi (come il \emph{weight decay}), garantendo che l'ottimizzatore minimizzi esclusivamente il puro errore quadratico medio (MSE):
\begin{equation}
\mathcal{L}_{\text{MSE}} = \frac{1}{N} \sum_{i=1}^N \lVert \mathbf{x}_i - g_{\theta}(f_{\phi}(\mathbf{x}_i)) \rVert^2
\end{equation}

Sotto tali condizioni, assumendo i dati in ingresso centrati, le equazioni della rete si riducono a una doppia trasformazione matriciale pura: $\mathbf{z} = \mathbf{W}_1 \mathbf{x}$ per l'Encoder e $\mathbf{\hat{x}} = \mathbf{W}_2 \mathbf{z}$ per il Decoder.

Come dimostrato formalmente dallo studio seminale di Bourlard e Kamp (1988), l'addestramento di un Linear AE così vincolato comporta una convergenza geometrica rigorosa verso l'ottimo globale. L'algoritmo di retropropagazione farà sì che lo spazio latente generato dalle matrici dei pesi $\mathbf{W}_1$ e $\mathbf{W}_2$ ricopra esattamente il medesimo sottospazio principale identificato dalle prime $k$ componenti della PCA. 

Sebbene le matrici dei pesi dell'AE non siano necessariamente ortogonali come nella PCA, l'Autoencoder Lineare funge da prova matematica cruciale: le reti neurali sono in grado di replicare perfettamente l'ottimo globale lineare. Questo lo rende il \emph{benchmark} strutturale di partenza prima di introdurre la complessità profonda.

\section{Autoencoder Profondi (AE): introduzione alla non-linearità}
\label{sec:deep_ae}

L'evoluzione naturale dell'architettura lineare è l'Autoencoder Profondo (Deep AE). Aggiungendo strati nascosti (\emph{hidden layers}) sovrapposti e introducendo funzioni di attivazione non lineari (come ReLU, Tanh o Sigmoide), la rete neurale acquisisce la capacità, sancita dal Teorema di Approssimazione Universale, di modellare qualsiasi funzione continua complessa.

In questo regime, l'Encoder e il Decoder cessano di essere semplici matrici di proiezione e si trasformano in operatori topologici non lineari. Il Deep AE è in grado di apprendere una \emph{manifold} curva e contorta, piegando lo spazio a latente per mappare in modo efficiente le relazioni gerarchiche settoriali e le proprietà non gaussiane delle matrici di correlazione. La dimensionalità $D$ originale viene progressivamente ridotta attraverso gli strati fino alla dimensione latente $k$.

Tuttavia, per quanto formidabile nella fase di compressione e ricostruzione storica, il Deep AE fallisce qualora venga impiegato come motore stocastico per simulazioni Monte Carlo. Il problema risiede nella natura non vincolata del suo spazio latente. La rete apprende a disporre i campioni storici nello spazio $\mathbf{z}$ unicamente per minimizzare l'errore di ricostruzione, ignorando l'organizzazione globale di tale spazio. Il risultato è un dominio latente altamente irregolare, frammentato, popolato da vaste "zone vuote" prive di significato geometrico. Se si estraesse un vettore $\mathbf{z}_{\text{random}}$ casuale da questo spazio e lo si facesse passare per il Decoder, la rete genererebbe una matrice di correlazione sintetica totalmente priva di senso economico, incapace di rispettare i fatti stilizzati (violazione del limite di Marchenko-Pastur o della proprietà di Perron-Frobenius). Il Deep AE è eccellente nel ricordare, ma strutturalmente inadatto a immaginare.

\section{Variational Autoencoders (VAE) e regolarizzazione dello spazio latente}
\label{sec:vae}

Per abilitare la generazione stocastica di scenari realistici, è necessario un cambio di paradigma architetturale che conduca dal determinismo geometrico alla teoria dell'inferenza probabilistica variazionale. Il \emph{Variational Autoencoder} (VAE), introdotto da Kingma e Welling (2013), risolve il problema della frammentazione dello spazio latente costringendo l'Encoder a emettere non più singoli vettori puntuali, ma intere distribuzioni di probabilità.

Nel VAE, per ogni campione di input $\mathbf{x}$, l'Encoder (ora definito come rete di inferenza) prevede due vettori distinti: un vettore delle medie $\boldsymbol{\mu}$ e un vettore delle varianze $\boldsymbol{\sigma}^2$. Questi vettori parametrizzano una distribuzione gaussiana multivariata $q_{\phi}(\mathbf{z}|\mathbf{x}) = \mathcal{N}(\boldsymbol{\mu}, \text{diag}(\boldsymbol{\sigma}^2))$ dalla quale viene campionato il vettore latente $\mathbf{z}$. Affinché questo processo di estrazione stocastica sia differenziabile tramite \emph{backpropagation}, si implementa il cosiddetto \emph{Reparameterization Trick}:
\begin{equation}
\mathbf{z} = \boldsymbol{\mu} + \boldsymbol{\sigma} \odot \boldsymbol{\epsilon} \quad \text{con} \quad \boldsymbol{\epsilon} \sim \mathcal{N}(\mathbf{0}, \mathbf{I})
\end{equation}

La grandezza ingegneristica del VAE risiede nella sua funzione di \emph{loss}, che massimizza il limite inferiore dell'evidenza (ELBO). Essa è composta da due termini antagonisti:
\begin{equation}
\mathcal{L}_{\text{VAE}} = \underbrace{\mathbb{E}_{q_{\phi}(\mathbf{z}|\mathbf{x})} [-\log p_{\theta}(\mathbf{x}|\mathbf{z})]}_{\text{Reconstruction Loss}} + \underbrace{\beta \, D_{KL}(q_{\phi}(\mathbf{z}|\mathbf{x}) \parallel p(\mathbf{z}))}_{\text{KL Divergence}}
\end{equation}

Il primo termine è la \emph{Reconstruction Loss} (tipicamente calcolata come MSE nel nostro dominio continuo), che spinge la rete a minimizzare l'errore tra la matrice di correlazione in ingresso e quella ricostruita. 

Il secondo termine, la Divergenza di Kullback-Leibler ($D_{KL}$), è il motore stocastico dell'architettura. Essa agisce come un rigoroso regolarizzatore topologico, penalizzando l'Encoder se le distribuzioni generate $q_{\phi}(\mathbf{z}|\mathbf{x})$ si allontanano da una distribuzione \emph{prior} di riferimento, solitamente una gaussiana normale standard $p(\mathbf{z}) = \mathcal{N}(\mathbf{0}, \mathbf{I})$.

La competizione tra questi due termini plasma uno spazio latente rivoluzionario. La Divergenza KL costringe tutte le distribuzioni dei singoli campioni ad addensarsi vicino all'origine dello spazio e a sovrapporsi le une con le altre in modo continuo e omogeneo (\emph{packing}), eliminando le zone vuote tipiche degli AE classici. Contemporaneamente, la Reconstruction Loss costringe i campioni simili (es. matrici di correlazione di fasi di espansione economica) a rimanere vicini per poter essere decodificati correttamente, preservando un'ordinata morfologia strutturale.

Questo spazio latente continuo, denso e regolarizzato è il dominio perfetto per il campionamento stocastico. Nelle fasi di simulazione di portafoglio, sarà sufficiente campionare un vettore casuale $\mathbf{z}_{\text{sim}} \sim \mathcal{N}(\mathbf{0}, \mathbf{I})$ e passarlo al Decoder deterministico $p_{\theta}(\mathbf{x}|\mathbf{z})$. Il modello estrarrà dalla regione di sovrapposizione un pattern coerente, generando una matrice di correlazione sintetica totalmente nuova, mai apparsa nella storia passata, ma che rispetterà rigorosamente l'infrastruttura topologica e i vincoli finanziari (fatti stilizzati) del mercato di addestramento.